% s-062-p062-m002; original PDF 62, 63
\Aphorism{s-062-p062-m002}{61}
Die älteren Philosophen wollten den Menschen abgehärtet und als ein geistiges Pachyderma den Berührungen der Umwelt ausgesetzt; in dieser Klugheit bestand oft genug ihre ganze Weisheit. Was wir Neueren erstreben, ist in ihrem Sinne gewiss unphilosophisch und keine geringe Donquixoterie: dem Individuum bei gegebenem Volumen ein Maximum von Oberfläche anbilden, seine Verwundbarkeit, die Fülle seiner Beziehungen, das Bereich seiner Erlebnisse steigern, seine Empfindlichkeit auf alle Töne und Farben und Vierteltöne und Zwischenfarben stimmen. Dieses Ideal mag jedem eudämonistisch gerichteten Willen widerstehen — und was war bisher nicht eudämonistisch? Die Moral, die ein beruhigtes Gewissen, und die Wissenschaft, die eine befestigte Gewissheit an die Spitze der Güter stellt, waren es in erster Reihe; überall herrscht die bleiche Furcht, die Scheu vor dem problematischen Charakter des Daseins, der Wunsch nach verminderter Unsicherheit, der Stoicismus, der um das Leben herum oder durch das Leben hindurch will, ohne seinen Stachel zu fühlen. Aber wer erkennen will, muss nackt ins Meer hinaus schwimmen und darf sich durch keine Hülle der vollen abenteuerlichen Willkür im Leben und Leiden entziehen.
\par

